\documentclass[UTF8,11pt]{ctexart}
\usepackage[a4paper,margin=22mm]{geometry}
\usepackage{graphicx,enumitem,amssymb,tabularx}
\usepackage[normalem]{ulem}
\usepackage{xeCJKfntef}
\setlength{\parindent}{0pt}
\setlength{\parskip}{.6em}
\sloppy
\begin{document}
\section*{2023年12月英语四级真题(第2套)}
2023 年12 月大学英语四级考试真题（第2 套）


\subsection*{\textbf{Part I Writing} \textbf{( 30 minutes)}}


\textbf{Directions:} \textit{Suppose the university newspaper is inviting submissions from the students for its coming edition on} \textit{the recent developments in their hometown. You are now to write an essay for submission. You will have 30} \textit{minutes to write the essay. You should write at least }\textit{\uline{120 }}\textit{words but no more than }\textit{\uline{180 }}\textit{words.}


\subsection*{\textbf{Part II} \textbf{Listening Comprehension} \textbf{( 25 minutes)}}


\subsection*{\textbf{Section A}}


\textbf{Directions:} \textit{In this section, you will hear three news reports. At the end of each news report, you will hear two or} \textit{three questions. Both the news report and the questions will be spoken only once. After you hear a question, you} \textit{must choose the best answer from the four choices marked A), B), C) and D). Then mark the corresponding letter} \textit{on }\textbf{\textit{Answer Sheet 1}}\textit{ with a single line through the centre.}


\textbf{Questions 1 and 2 are based on the news report you have just heard.}


1.


\begin{itemize}[leftmargin=2em,itemsep=.2em,topsep=.4em]\item[A.] Their interests are quite similar.
\item[B.] They are generally the same age.
\item[C.] Their brains work in harmony.
\item[D.] They have the same ethnic background.
\end{itemize}

\noindent\begin{tabularx}{\linewidth}{@{}lXXXX@{}}2. & \textbf{A.} It can work both ways. & \textbf{B.} It can be touching. & \textbf{C.} It is hard to predict. & \textbf{D.} It resembles family ties.\\\end{tabularx}\par

\textbf{Questions 3 and 4 are based on the news report you have just heard.}


3.


\begin{itemize}[leftmargin=2em,itemsep=.2em,topsep=.4em]\item[A.] Identify their biological fathers.
\item[B.] Search for their half-brothers.
\item[C.] See whether they are actually related.
\item[D.] Find out more about their ancestry.
\end{itemize}

4.


\begin{itemize}[leftmargin=2em,itemsep=.2em,topsep=.4em]\item[A.] They were both 60 years of age.
\item[B.] They were born to the same mother.
\item[C.] They flew 737 airplanes as pilots.
\item[D.] They were both given up for adoption.
\end{itemize}

\textbf{Questions 5 to 7 are based on the news report you have just heard.}


5.


\begin{itemize}[leftmargin=2em,itemsep=.2em,topsep=.4em]\item[A.] The kilometer-long beach was practically deserted.
\item[B.] The beautiful beach was spoiled with lots of trash
\item[C.] Other tourists refused to join in the cleanup.
\item[D.] One of his friends was caught littering.
\end{itemize}

\noindent\begin{tabularx}{\linewidth}{@{}lXXXX@{}}6. & \textbf{A.} The beach authorities. & \textbf{B.} One of the five tourists. & \textbf{C.} A passerby. & \textbf{D.} A local woman.\\\end{tabularx}\par

7.


\begin{itemize}[leftmargin=2em,itemsep=.2em,topsep=.4em]\item[A.] It was tourists not natives who were cleaning up the beach.
\item[B.] The number of tourists to the beach is on a steady decline.
\item[C.] Some natives were selling poor-quality food to tourists.
\item[D.] The tourists’ good deed was not noticed by the locals.
\end{itemize}

\subsection*{\textbf{Section B}}


\textbf{Directions:} \textit{In this section, you will hear two long conversations. At the end of each conversation, you will hear} \textit{four questions. Both the conversation and the questions will be spoken only once. After you hear a question, you} \textit{must choose the best answer from the four choices marked A), B), C) and D). Then mark the corresponding letter} \textit{on }\textbf{\textit{Answer Sheet 1}}\textit{ with a single line through the centre.}


\textbf{Questions 8 to 11 are based on the conversation you have just heard.}


8.


\begin{itemize}[leftmargin=2em,itemsep=.2em,topsep=.4em]\item[A.] He has to play football with workmates.
\item[B.] He has got some books to read.
\item[C.] He is going to visit a friend.
\item[D.] He is physically unfit for it.
\end{itemize}

9.


\begin{itemize}[leftmargin=2em,itemsep=.2em,topsep=.4em]\item[A.] To teach kids about animal protection.
\item[B.] To learn how popular zoos could be.
\item[C.] To see some rare animals in cages.
\item[D.] To give her little nephew a treat.
\end{itemize}

10.


\begin{itemize}[leftmargin=2em,itemsep=.2em,topsep=.4em]\item[A.] He enjoys excellent health.
\item[B.] He is keen on extreme sports.
\item[C.] He coaches tennis players every week.
\item[D.] He spends most of his time in the gym.
\end{itemize}

11.


\begin{itemize}[leftmargin=2em,itemsep=.2em,topsep=.4em]\item[A.] Tending to his swollen ankle.
\item[B.] Concentrating on reading.
\item[C.] Writing three book reports.
\item[D.] Planning Christmas celebrations.
\end{itemize}

\textbf{Questions 12 to 15 are based on the conversation you have just heard.}


12.


\begin{itemize}[leftmargin=2em,itemsep=.2em,topsep=.4em]\item[A.] It is being debated by hundreds of retirees.
\item[B.] It is attracting many people’s attention.
\item[C.] It partly records his own experience.
\item[D.] It argues for postponing retirement.
\end{itemize}

13.


\begin{itemize}[leftmargin=2em,itemsep=.2em,topsep=.4em]\item[A.] One should foresee a financial crisis.
\item[B.] One should trust financial planners’ figures.
\item[C.] One should have one million dollars to retire.
\item[D.] One should start saving as early as possible.
\end{itemize}

14.


\begin{itemize}[leftmargin=2em,itemsep=.2em,topsep=.4em]\item[A.] It doesn’t need to be permanent.
\item[B.] It shouldn’t be considered risky.
\item[C.] It helps to reduce travel expenses.
\item[D.] It is the way to quit a job one hates.
\end{itemize}

15.


\begin{itemize}[leftmargin=2em,itemsep=.2em,topsep=.4em]\item[A.] By keeping close contact with one’s employers.
\item[B.] By following the counsel of financial planners.
\item[C.] By retiring when one reaches sixty years old.
\item[D.] By investing half of one’s monthly income.
\end{itemize}

\subsection*{\textbf{Section C}}


\textbf{Directions:}\textit{ In this section, you will hear three passages. At the end of each passage, you will hear three or four} \textit{questions. Both the passage and the questions will be spoken only once. After you hear a question, you must} \textit{choose the best answer from the four choices marked A), B), C) and D). Then mark the corresponding letter on} \textbf{\textit{Answer Sheet 1}}\textit{ with a single line through the centre.}


\textbf{Questions 16 to 18 are based on the passage you have just heard.}


16.


\begin{itemize}[leftmargin=2em,itemsep=.2em,topsep=.4em]\item[A.] They tended to be arbitrarily judged by individuals of opposing groups.
\item[B.] They tended to be easily anticipated by those belonging to their own race.
\item[C.] They were influenced by the presence of someone from an outsider group.
\item[D.] They were readily shared among members of the same social or racial group.
\end{itemize}

17.


\begin{itemize}[leftmargin=2em,itemsep=.2em,topsep=.4em]\item[A.] When an unknown student from another university was present
\item[B.] When an experimenter from the research team took notice.
\item[C.] When they were offered both candy and fruit as a snack.
\item[D.] When they tried to make a positive impression on the researchers.
\end{itemize}

18.


\begin{itemize}[leftmargin=2em,itemsep=.2em,topsep=.4em]\item[A.] By maintaining its positive image.
\item[B.] By advertising its social benefits.
\item[C.] By supporting struggling consumers.
\item[D.] By teaching consumers diet strategies.
\end{itemize}

\textbf{Questions 19 to 21 are based on the passage you have just heard.}


19.


\begin{itemize}[leftmargin=2em,itemsep=.2em,topsep=.4em]\item[A.] The academic and learning issues struggling students encounter.
\item[B.] The risk students face due to a history of mental health problems.
\item[C.] The work universities are doing to help students succeed academically.
\item[D.] The effect of interacting with therapy dogs on students under pressure.
\end{itemize}

20.


\begin{itemize}[leftmargin=2em,itemsep=.2em,topsep=.4em]\item[A.] Their executive functioning.
\item[B.] Their communicative skills.
\item[C.] Their academic networking.
\item[D.] Their leadership capacities.
\end{itemize}

21.


\begin{itemize}[leftmargin=2em,itemsep=.2em,topsep=.4em]\item[A.] Rid students of their anxiety.
\item[B.] Add to some students’ stress.
\item[C.] Contribute little to typical students’ success.
\item[D.] Help students with mental issues pull through.
\end{itemize}

\textbf{Questions 22 to 25 are based on the passage you have just heard.}


22.


\begin{itemize}[leftmargin=2em,itemsep=.2em,topsep=.4em]\item[A.] Work hard and plan carefully.
\item[B.] Attempt to succeed at any cost.
\item[C.] Aim high and expect great results.
\item[D.] Remain optimistic even in difficulty.
\end{itemize}

23.


\begin{itemize}[leftmargin=2em,itemsep=.2em,topsep=.4em]\item[A.] Regarding failure as something inevitable.
\item[B.] Trying out innovative marketing strategies.
\item[C.] Venturing into sectors never explored before.
\item[D.] Being willing to experiment with novel ideas.
\end{itemize}

24.


\begin{itemize}[leftmargin=2em,itemsep=.2em,topsep=.4em]\item[A.] Expect future success so as to move forward.
\item[B.] Learn from our failure and forge ahead.
\item[C.] Distinguish between good and bad risks.
\item[D.] Examine our strategies and find out weaknesses.
\end{itemize}

\noindent\begin{tabularx}{\linewidth}{@{}lXXXX@{}}25. & \textbf{A.} Fresher offers. & \textbf{B.} Safer operation. & \textbf{C.} More challenges. & \textbf{D.} Less competition.\\\end{tabularx}\par

\subsection*{\textbf{Part III Reading Comprehension ( 40 minutes)}}


\subsection*{\textbf{Section A}}


\textbf{Directions:} \textit{In this section, there is a passage with ten blanks. You are required to select one word for each blank} \textit{from a list of choices given in a word bank following the passage. Read the passage through carefully before} \textit{making your choices. Each choice in the bank is identified by a letter. Please mark the corresponding letter for} \textit{each item on }\textbf{\textit{Answer Sheet 2}}\textit{ with a single line through the centre. You may not use any of the words in the bank} \textit{more than once.}


A number of studies have looked at how family life can affect productivity and satisfaction in the workplace. However, there has been \uline{26} little research on the influence of leisure activities. So Ciara Kelly and colleagues recruited 129 hobbyists to look at how the time spent on their hobbies \uline{27} their work life.


The researchers found that when participants spent longer than \uline{28} on their leisure activity, their belief in their ability to perform their job was strengthened. But this was only the \uline{29} when they had a serious hobby that was dissimilar to their job, or when their hobby was similar to their work but they only did it \uline{30 }. When their hobby was both serious and similar to their job, then spending more time on it actually decreased their work \uline{31 }.


Why might that be? To maintain a serious hobby, people need to invest significant psychological resources, say the authors — so if the activity has the same kinds of demands as their work, they may be left \uline{32} and unable to perform well at their job. But if their hobby is quite different from their career, it may not \uline{33} in the same way but instead help them develop other knowledge and skills that can \uline{34} their confidence at work. “Consider a scientist who is a keen rock climber,” says Kelly. “Since climbing is so far


\uline{35} from their day-to-day work activities, they can still recover from the demands of their job with plenty of resources.”


\begin{itemize}[leftmargin=2em,itemsep=.2em,topsep=.4em]\item[A.] boost
\item[B.] case
\item[C.] casually
\item[D.] efficiency
\item[E.] estate
\item[F.] exhausted
\item[G.] faculty
\item[H.] interfere
\item[I.] normal
\item[J.] prevalent
\item[K.] relative
\item[L.] removed
\item[M.] scratch
\item[N.] shaped
\item[O.] surprisingly
\end{itemize}

\subsection*{\textbf{Section B}}


\textbf{Directions:} \textit{In this section, you are going to read a passage with ten statements attached to it. Each statement} \textit{contains information given in one of the paragraphs. Identify the paragraph from which the information is derived.} \textit{You may choose a paragraph more than once. Each paragraph is marked with a letter. Answer the questions by} \textit{marking the corresponding letter on }\textbf{\textit{Answer Sheet 2}}\textit{.}


\subsection*{\textbf{More fathers are taking paternity leave, but mothers are still doing all the work}}


A) Attitudes towards paternity leave (陪产假) have drastically changed in America in the last five years as more fathers feel comfortable taking extended time off, but gender bias persists when it comes to career prospects and the home, according to a new study of working parents.


B) Research by the Boston College Center for Work \& Family, which surveyed new parents at four large U.S. companies who were qualified for taking at least six weeks paid parental leave, found that 81\% of the 1,240 employees surveyed said the notion of fathers taking leave has become more acceptable.


C) Of those surveyed, 62\% of fathers took the maximum amount of time off compared to 93\% of mothers, and around three-quarters of workers said their employer was equally supportive of mothers and fathers taking parental leave and over half said leave policies had made workplace culture better.


D) The U.S. is one of only three countries in the world not to offer statutory (法定的) paid leave, but increasingly states and companies are starting to take up the issue. So far, eight states and the District of Columbia have their own paid family leave laws.


E) Brad Harrington, executive director of the center and lead author of the study, estimates only 20\% to 30\% of companies in the U.S. offer paid parental leave. He feels the research findings reflect an obvious change in corporate attitudes to new fathers taking time off.


F) “We did a study on paternity leave five years ago. Compared with those findings, these numbers were shocking to me. I did not expect 80\% of people to say the organisation finds dads taking this leave acceptable and three-quarters to say it’s equally supportive of women and men taking leave,” Harrington said.


G) The previous study found that nearly three-quarters of fathers saw two to four weeks as an appropriate duration for paternity leave and 76\% said they would prefer not to take all their time off at once.


H) Since then, there have been a number of legal cases against companies involving paternity leave — including cases against JPMorgan Chase and Estée Lauder — which have helped put pressure on companies to make their parental leave policies gender neutral.


I) However, the study also shows how traditional gender roles endure both at work, where more women than men reported changes in their perceived career potential, and at home — even among workers who claim to have a strong desire for equality.


J) The vast majority of men, 97\%, said one of the top reasons to take leave was to share caregiving with their partner. But when they were asked about how caregiving and household tasks were divided, their answers painted a different picture. While about 75\% of employees said both genders should give the same amount of care, the majority of men and nearly half of women admitted that in reality the female actually did most of it. A tiny fraction, 2\%, of men said they did more of the childcare.


K) Men and women’s experiences of the return to work following parental leave were also considerably different. Of the women surveyed, 32\% reported a downturn in their job satisfaction, while 14\% said it increased. In comparison, 17\% of men said their job satisfaction went down and 20\% said it went up. Meanwhile, more women reported an increase in their responsibilities and manager expectations after childbirth. Half of the women said they used flexible work arrangements after becoming a parent, while just 27\% of men did. Similar percentages of men and women said they enjoyed their careers and that it gave them a sense of achievement, while around half of women and 44\% of men said it was a key part of their identity.


L) On the subject of career advancement, 59\% of women and 49\% of men said leave could be limiting and both genders said they feared it would have an impact on their progress long-term. But on opportunity for promotion, more than double the number of women, 30\% compared with 15\% of men, believed their chances to be lower after becoming a parent. Despite progress, the struggle for women to reach the highest positions of power is demonstrated in this year’s Fortune 500 list, which featured a record 33 female CEOs, but this still represents a tiny fraction of the total.


M) Harrington said culture change depends on companies putting more focus on men and their responsibilities. “By that I mean companies need to give men paternity leave and encourage men to take time off to be with their kids early on in the kids’ life. They also need to recognise that men have to make significant adjustments when they become parents. Companies cannot do all these things to enhance women’s advancement and then turn around and say, ‘Oh, but we don’t expect the men to take over for the women at home.’ ”


N) In May, the American Civil Liberties Union (ACLU) and Outten \& Golden LLP announced a historic class-action \$5m settlement with JPMorgan Chase on behalf of male employees who claim they were illegally denied access to paid parental leave. Derek Rotondo, 35, filed the discrimination charge against his company after he was allegedly told by his HR department that mothers were considered primary caregivers. Thus, they were allowed to take 16 weeks of paid parental leave. Fathers, however, could take just two weeks.


O) The father of two from Columbus, Ohio, who still works at the company as an associate and investigator, said he has witnessed a “domino effect (多米诺效应)” across companies since the settlement, but that there is still substantial progress to be made towards changing attitudes towards paternity leave.


P. “I do think there’s still some way to go...there’s still going to be sort of the unstated expectation for new dads to essentially come right back to work, but I think the research is showing that’s starting to change.” He said equal parental leave is an essential component to creating gender equality in the workplace. “The old standard of women staying home, having babies and cooking doesn’t apply and hasn’t applied for a long time.”


36. In the absence of Federal legislation, some states in the U.S. have passed laws concerning paid family leave.


37. Most fathers admitted that even during their paternity leave they actually did much less childcare than the


mother.


38. According to one father, equal parental leave is indispensable to achieving gender equality in the workplace.


39. One survey indicated there is now less objection to paternity leave.


40. Compared to five years ago, according to one researcher, many more people said their organisation gave the same support to men and women taking parental leave.


41. One study finds that even workers who claim to desire gender equality stick to traditional gender roles both at work and at home.


42. The majority of workers surveyed said parental leave policies had improved workplace culture.


43. In spite of progress, the number of women in top positions of big companies remains extremely small.


44. According to one estimate, less than one third of companies in the U.S. provide paid parental leave.


45. A number of lawsuits have pressured companies to formulate gender neutral policies on parental leave.


\subsection*{\textbf{Section C}}


\textbf{Directions: }\textit{There are 2 passages in this section. Each passage is followed by some questions or unfinished} \textit{statements. For each of them there are four choices marked A), B), C) and D). You should decide on the best} \textit{choice and mark the corresponding letter on }\textbf{\textit{Answer Sheet 2}}\textit{ with a single line through the centre.}


\subsection*{\textbf{Passage One}}


\textbf{Questions 46 to 50 are based on the following passage.}


Having a rival can keep you committed to achieving your goals and enhance your overall performance. But before you go out and find an entrepreneur to outcompete, it’s important to understand and avoid the traps that often come with rivalry. After all, competitive rivalry can also hinder effective decision-making and increase your willingness to take risks, behaviors that can ultimately hurt your venture’s success.


Finding someone you’re committed to outcompeting can be a great way to stay focused on your goals and push your venture to the next level. But when you’re intently focused on outperforming your rivals, you may begin to develop a “win-at-all-costs” mentality that causes you to ignore how you achieve success. One group of researchers, for example, examined the link between rivalry and unethical behavior. They found that when people compete against their rivals, they are more willing to behave unethically to win. But such behavior may stain your reputation and strain relationships important to your success. One way to avoid this trap is to stop and reflect on what’s important. While outperforming your rivals may provide short-term benefits, the loss of your integrity will have long-term consequences.


One reason having a rival can enhance your venture’s performance is that it creates a level of excitement that drives you to work harder. But this eagerness to win may also hurt your venture’s success, particularly when it causes you to make impulsive, insensible decisions. But it’s possible to avoid such costly mistakes by making a habit of engaging in critical thinking, such as considering opposing viewpoints and conducting cost-benefit analyses, especially for those decisions that are complex and can determine the future of your venture.


The sense of eagerness that comes with having a rival can not only cause you to make poorer decisions, but it can also lead you to take greater risks that put your venture in peril. One way you can overcome the risk-inducing effects of rivalry that stand to endanger your venture’s success is to remain attentive to your emotional state and actively monitor how such feelings are affecting your decision-making.


46. How can competitive rivalry benefit entrepreneurs according to the passage?


\begin{itemize}[leftmargin=2em,itemsep=.2em,topsep=.4em]\item[A.] By enabling them to outcompete other entrepreneurs.
\item[B.] By enabling them to make their venture a success.
\item[C.] By helping them to reach long-term goals.
\item[D.] By helping them to stay goal-oriented.
\end{itemize}

47. What is one of the traps entrepreneurs may often fall into when competing with rivals?


\begin{itemize}[leftmargin=2em,itemsep=.2em,topsep=.4em]\item[A.] They may adopt strategies that are bound to ruin their venture.
\item[B.] They may resort to unethical means to outperform their rivals.
\item[C.] They may be too eager to succeed while ignoring the huge labor cost.
\item[D.] They may be intently focused on winning at the current market level.
\end{itemize}

48. What are entrepreneurs advised to do to avoid traps that often accompany rivalry?


\begin{itemize}[leftmargin=2em,itemsep=.2em,topsep=.4em]\item[A.] Deliberate on what really matters.
\item[B.] Prioritize reaping immediate benefits.
\item[C.] Estimate the long-term consequences.
\item[D.] Reflect on what successes are achievable.
\end{itemize}

49. How can entrepreneurs avoid making impulsive and insensible decisions?


\begin{itemize}[leftmargin=2em,itemsep=.2em,topsep=.4em]\item[A.] By engaging themselves in critical reasoning.
\item[B.] By developing a habit of keeping their integrity.
\item[C.] By criticizing themselves for previous poor performances.
\item[D.] By refraining from being too excited about their successes.
\end{itemize}

50. How can entrepreneurs overcome the risk-inducing effects of rivalry?


\begin{itemize}[leftmargin=2em,itemsep=.2em,topsep=.4em]\item[A.] By paying close attention to their current performance.
\item[B.] By taking steps that stand to endanger their rivals’ success.
\item[C.] By monitoring how their decision-making impacts their mentality.
\item[D.] By keeping their emotions in check to avoid making poor decisions.
\end{itemize}

\subsection*{\textbf{Passage Two}}


\textbf{Questions 51 to 55 are based on the following passage.}


A multitasker is one who can perform two or more tasks effectively at the same time, which — apart from the obvious differences — is similar to what a computer does. The concept does indeed come from the realms of technology, where it is used to refer to an operating system that can execute multiple tasks at the same time. However, the question is: can a person really be a multitasker?


For most scientists, the answer is no. So much so that, according to experts in neuroscience(神经系统科学), our brains do not handle multitasking situations well. As soon as two tasks require our attention, productivity suffers. What we call multitasking, therefore, is in reality the ability to move more or less quickly from one task to another. This requires two essential conditions: that one of the tasks needs to be automatic, like walking or eating, and that they both need different mental processes. Answering the phone and writing at the same time, for example.


However, on the other side of the coin there are people who maintain that it is possible to be, or at least seem to be, multitasking. A recent study concluded that regardless of whether people are actually handling several tasks or not, the mere fact that they perceive this activity as multitasking has a positive effect on their performance.


The business perspective offers a different view: multitasking is understood as the ability to adapt to all types of environment within a company and effectively undertake different activities within a set time frame. Indeed, many companies look for people who are skilled in multitasking to improve their productivity. From this different perspective, you can not only be multitasking but this ability can also be taught: something that is easier in fluid organisations, which favour flexibility in their working practices.


The benefits of multitasking are clear. Being quicker and more efficient increases our performance and the number of tasks completed. But having to pay attention to several things at once means that the powers of concentration are reduced and that can lead to more mistakes.


51. What does a “multitasker” originally refer to?


\begin{itemize}[leftmargin=2em,itemsep=.2em,topsep=.4em]\item[A.] An operating system capable of doing several tasks at once.
\item[B.] A skilled worker executing more than one task at the same time.
\item[C.] A sophisticated technology doing several tasks effectively at once.
\item[D.] An efficient person able to perform multiple tasks at the same time.
\end{itemize}

52. Why can’t people really be multitaskers according to neuroscientists?


\begin{itemize}[leftmargin=2em,itemsep=.2em,topsep=.4em]\item[A.] They are not sufficiently exposed to multitasking situations.
\item[B.] They are not comparable to mechanical operating systems.
\item[C.] Their brains do not allow them to multitask.
\item[D.] Their attention span cannot be expanded.
\end{itemize}

53. What do we learn from the conclusion of a recent study on multitasking?


\begin{itemize}[leftmargin=2em,itemsep=.2em,topsep=.4em]\item[A.] People make greater achievements by maintaining whatever they are doing is multitasking.
\item[B.] People’s performance benefits from the perception of what they are doing as multitasking.
\item[C.] People’s active mental processes exert a positive effect on their multitasking.
\item[D.] People can improve their capabilities by handling multitasking situations.
\end{itemize}

54. How does the business world view multitasking?


\begin{itemize}[leftmargin=2em,itemsep=.2em,topsep=.4em]\item[A.] It is a rare skill often found in fluid organisations.
\item[B.] It is an adaptable capability required of all workers.
\item[C.] It is an essential quality many employees lack.
\item[D.] It is a desirable ability that can be developed.
\end{itemize}

55. What does the author imply we should do if we have to focus on some task and do it well?


\begin{itemize}[leftmargin=2em,itemsep=.2em,topsep=.4em]\item[A.] Work in a flexible way.
\item[B.] Learn from mistakes.
\item[C.] Avoid multitasking.
\item[D.] Increase efficiency.
\end{itemize}

\subsection*{\textbf{Part IV Translation} \textbf{( 30 minutes)}}


\textbf{Directions: }\textit{For this part, you are allowed 30 minutes to translate a passage from Chinese into English. You} \textit{should write your answer on }\textbf{\textit{Answer Sheet 2}}\textit{.}


改革开放以来，中国人民生活水平不断提高,这在人们的\CJKunderline{饮食 }(diet) 变化上得到充分体现。如今，人们不再满足于吃得饱，而是追求吃得更加安全、更加营养、更加健康，食物也愈来愈丰富多样，不再限于本地的农产品。\CJKunderline{物流业 }(logistics industry) 的发展使人们很容易品尝到全国各地的特产。毫无疑问，食品质量与饮食结构的改善为增进人们健康提供了有力的保障。

\end{document}